% Generated draft: requires mathematical and source review.
\documentclass[11pt]{article}
\usepackage[margin=25mm]{geometry}
\usepackage{fontspec}
\setmainfont{Latin Modern Roman}
\usepackage{amsmath,amssymb,mathtools,unicode-math}
\setmathfont{Latin Modern Math}
\usepackage{enumitem,needspace,xurl,hyperref}
\setlength{\parindent}{0pt}
\setlength{\parskip}{5pt}
\setlength{\emergencystretch}{3em}
\newcommand{\N}{\mathbb{N}}
\newcommand{\Z}{\mathbb{Z}}
\newcommand{\Q}{\mathbb{Q}}
\newcommand{\R}{\mathbb{R}}
\newcommand{\C}{\mathbb{C}}
\newcommand{\F}{\mathbb{F}}
\newcommand{\A}{\mathbb{A}}
\newcommand{\PP}{\mathbb{P}}
\newcommand{\E}{\mathbb{E}}
\newcommand{\eps}{\varepsilon}
\newcommand{\dd}{\,\mathrm{d}}
\begin{document}
\section*{UnsolvedMath ID: 30003144}
\textbf{Source problem code:} catalogue entry
\subsection{Definitions and mathematical statement}
An even integral lattice is a finite-rank free abelian group equipped with a nondegenerate integral symmetric bilinear form \((\cdot,\cdot)\) such that \((x,x)\in 2\Z\) for every lattice vector \(x\). Use the hyperbolic sign convention \(\operatorname{sig}(S)=(n,1)\), so that a square-\(2\) vector means a vector \(r\in S\) with \((r,r)=2\). Its reflection is
\[
 s_r(x)=x-(x,r)r.
\]
The group generated by these reflections is the \(2\)-reflection Weyl group \(W^{(2)}(S)\). The supplied description requires this group to have finite index in the relevant orthogonal group and a fundamental chamber to admit a lattice Weyl vector. With chamber walls oriented by inward inequalities \((x,\alpha)\leq 0\), the Weyl-vector equations for square-\(2\) wall roots are \((\rho,\alpha)=-1\). The record does not specify the precise ambient orthogonal group or whether the Weyl vector must lie in \(S\) or in \(S\otimes\Q\).

Let \(\mathcal S\) denote the particular collection of \(59\) even hyperbolic lattices referred to in Theorem~3 of the supplied report. The supplied reconstruction describes them by the preceding finite-index and Weyl-vector properties. Neither their list nor the full hypotheses of that theorem are included, so \(\mathcal S\) remains a source-dependent collection; the description alone is not asserted here to be an exhaustive classification without further hypotheses.

Write \(U(m)\) for the rank-two lattice with Gram matrix
\[
 \begin{pmatrix}0&m\\m&0\end{pmatrix},
 \qquad T=U(m)\oplus S.
\]
For nonzero \(m\), the lattice \(T\) has signature \((n+1,2)\). Let \(\Omega(T)\) be a connected component of its type-IV domain
\[
 \left\{[z]\in\PP(T\otimes\C):(z,z)=0,\ (z,\overline z)<0\right\}.
\]
For \(r\in T\) with \((r,r)=2\), set
\[
 \mathcal D_r=\{[z]\in\Omega(T):(z,r)=0\},
 \qquad
 \mathscr H_2(T)=\{\mathcal D_r:r\in T,\ (r,r)=2\},
\]
where coincident hypersurfaces are counted only once, in particular \(\mathcal D_r=\mathcal D_{-r}\).

\textbf{Question.} For every \(S\in\mathcal S\), does there exist an integer \(m\) for which this construction is nondegenerate, and a nonzero holomorphic automorphic form \(\Phi\) on \(\Omega(T)\), of positive weight and with integral Fourier coefficients, such that
\[
 \operatorname{div}_{\Omega(T)}(\Phi)
   =\sum_{H\in\mathscr H_2(T)} H?
\]
Thus every square-\(2\) rational quadratic hypersurface must occur with multiplicity one, and there must be no other divisor components in the domain. This is the requested \(2\)-reflective form.

The integer \(m\), the weight, and \(\Phi\) may depend on \(S\). Nondegeneracy requires \(m\ne0\); the supplied record does not explicitly impose \(m>0\). The automorphy group, permitted character and weights, Fourier expansion conventions and normalization, and holomorphy requirements at the boundary are not specified in the record. Consequently these parts of the existence question require clarification from the cited source before the draft is a fully determined formulation.
\subsection{Short English statement}
For each of the report's specified \(59\) hyperbolic lattices, ask whether adjoining a scaled hyperbolic plane yields a domain supporting a positive-weight holomorphic automorphic form with integral Fourier coefficients, vanishing exactly once along every square-\(2\) reflection hypersurface. The supplied curation labels the problem partially solved and reports constructions for \(36\) of the \(59\) lattices; this dated assessment is not independently verified here. The lattice list and several automorphic conventions remain source-dependent.
\subsection{Sources}
\begin{itemize}
\item Local source registry: \texttt{data/unsolvedmath/problems.json}, with \texttt{catalog/problems.json} as the additional-record registry; supplied complete record for permanent UnsolvedMath ID \(30003144\), source problem code \texttt{catalogue entry}, entitled \emph{Two-Reflective Automorphic Forms for Hyperbolic Lattices}. No individual proposer is supplied; the proposed year is recorded as \(2016\).
\item Supplied report citation: \emph{Moduli spaces and Modular forms (2017)}, Oberwolfach Reports, pp.~1259--1318, DOI \texttt{10.4171/owr/2016/23}. \url{https://doi.org/10.4171/owr/2016/23}. The original statement refers to the \(59\) lattices in Theorem~3. The supplied verification note says that the official report was checked and that the reconstructed statement expands both that reference and the extended lattice. This is a supplied verification claim, not a claim of independent consultation in this draft.
\item Supplied literature source: Gritsenko--Nikulin, \emph{2-reflective Weyl groups}. \url{https://doi.org/10.1112/plms.12084}.
\item Supplied literature source, labeled ``primary paper, explicit 36-of-59 count'': \url{https://livrepository.liverpool.ac.uk/3019925/1/1602.08359v2.pdf}.
\item Supplied literature source, labeled ``current classification of 2-reflective modular forms'': \url{https://ems.press/content/serial-article-files/42563}.
\item Literature assessment checked at \texttt{2026-08-21}: the supplied curation states that the cited program constructs automorphic corrections for \(36\) of the \(59\) lattices and that the all-\(59\) existence assertion remains conjectural. This is supplied dated curation derived from the Oberwolfach Reports record, not an independently verified literature assessment. The record labels the status \texttt{partially\_solved} and the statement \texttt{corrected\_verified}; those labels are preserved as source metadata rather than established mathematical conclusions.
\end{itemize}
\end{document}
